\documentclass[10pt,a4paper]{amsart}
\usepackage[margin=19mm]{geometry}
\usepackage{amsmath,amssymb,mathtools}
\usepackage[scr=rsfs,cal=euler]{mathalfa}
\usepackage{xfrac}
\usepackage[unicode,hidelinks,bookmarks=false]{hyperref}
\newcommand{\ie}{i.e.\ }
\newcommand{\cf}{cf.\ }
\newcommand{\resp}{respectively\ }
\newcommand{\Subsection}{\subsection}
\providecommand{\nobreakdash}{\nobreak\hbox{-}}
\setlength{\emergencystretch}{2em}
\multlinegap0pt
\usepackage{amssymb}
\usepackage{mathtools}
\usepackage{tikz}
\usepackage[shortlabels]{enumitem}
\usepackage{xcolor}
\newtheorem{cor}{Corollary}
\newtheorem{lem}{Lemma}
\newtheorem{thm}{Theorem}
\usepackage{cleveref}
\crefname{cor}{Corollary}{Corollarys}
\crefname{lem}{Lemma}{Lemmas}
\crefname{thm}{Theorem}{Theorems}
\title{Pseudo-Gorenstein$^{*}$ Graphs}
\author{Takayuki Hibi, Selvi Kara, Dalena Vien}
\date{Source: arXiv:2603.08502v1 (CC BY 4.0)}
\begin{document}
\maketitle

\noindent\emph{Source and licence.} Pseudo-Gorenstein$^{*}$ Graphs. Version arXiv:2603.08502v1 (CC BY 4.0), distributed by arXiv under the Creative Commons Attribution 4.0 International licence, http://creativecommons.org/licenses/by/4.0/, as recorded in the arXiv licence metadata for this exact version and shown on the abstract page https://arxiv.org/abs/2603.08502v1. Full licence notice: \texttt{licenses/arxiv-2603.08502-CC-BY-4.0.txt}.\par\smallskip

\noindent\textit{Editorial scope.} Counted unit: the article's thm thm:PGstar-full-susp-path-cycle (source0.tex lines 790--799). The article's complete original proof of it is delivered with it (source0.tex lines 801--864, one \texttt{proof} environment of 64 lines). No mathematics is written, repaired or re-derived: the statement and the proof are the article's own lines, in the article's own notation, and no retained line is substituted or re-flowed.  Retained uncounted as the source-local context the proof needs: lines 296--298; lines 343--357; lines 378--394.  Nothing was omitted from any retained span, so the package carries no omission record.  No retained line contains a citation, so the wrapper carries no bibliography and no bibliography entry.  No retained line contains a figure environment, an image-inclusion command or drawing code, so no image is delivered and the package declares no asset dependency.  Editorial formatting only, in addition: an AMS \texttt{amsart} preamble that restates the article's own declarations and macros used by the retained lines, namely the environments \texttt{cor}, \texttt{lem}, \texttt{thm} on separate counters, together with the standard packages those lines need.  The retained environments print in this document as Corollary 1, Lemma 1, Theorem 1 in order of appearance, whereas the article prints the same items with the numbering of its own full document.  The complete native source of the article as distributed in the arXiv e-print for this exact version is bundled unchanged as \texttt{sources/196010/source0.tex}.
\begin{cor}\label{cor:PGstar-criterion}
A finite simple graph \(G\) is pseudo-Gorenstein\(^{*}\) if and only if \(P_G(-1)=(-1)^{\alpha(G)}\).
\end{cor}

\begin{lem}\label{lem:path-minus-one}
Let \(p_n:=P_{P_n}(-1)\) for \(n\ge 0\). Then
\[
p_0=1,\qquad p_1=0,\qquad p_n=p_{n-1}-p_{n-2}\ \ \text{for } n\ge 2.
\]
Consequently, \(p_n\) is periodic of period \(6\), and
\[
p_n=
\begin{cases}
1,& n\equiv 0,5 \pmod 6,\\[4pt]
0,& n\equiv 1,4 \pmod 6,\\[4pt]
-1,& n\equiv 2,3 \pmod 6.
\end{cases}
\]
\end{lem}

\begin{lem}\label{lem:cycle-minus-one}
Let \(c_n:=P_{C_n}(-1)\) for \(n\ge 3\). Then
\[
c_n=p_{n-1}-p_{n-3},
\]
where \(p_n\) is as in \Cref{lem:path-minus-one}. In particular, \(c_n\) is
periodic of period \(6\), and
\[
c_n=
\begin{cases}
2,& n\equiv 0 \pmod 6,\\[4pt]
1,& n\equiv 1,5 \pmod 6,\\[4pt]
-1,& n\equiv 2,4 \pmod 6,\\[4pt]
-2,& n\equiv 3 \pmod 6.
\end{cases}
\]
\end{lem}

\begin{thm}\label{thm:PGstar-full-susp-path-cycle}
Let \(\widehat{G}\) denote the \emph{full suspension} (cone) over a graph \(G\).
\begin{enumerate}
\item[(a)] For \(n\ge 3\), the full suspension \(\widehat{C_n}\) is pseudo-Gorenstein$^{*}$
if and only if \(n\equiv 0 \pmod{12}\).

\item[(b)] For \(n\ge 1\), the full suspension \(\widehat{P_n}\) is pseudo-Gorenstein$^{*}$
if and only if \(n\equiv 1,10 \pmod{12}\).
\end{enumerate}
\end{thm}

\begin{proof}
Let \(z\) be the new vertex in the full suspension \(\widehat{G}\).
Since \(z\) is adjacent to every vertex of \(G\), we have \(P_{\widehat{G}}(x)=P_G(x)+x\). Hence
\begin{equation}\label{eq:cone-eval-minus-one}
P_{\widehat{G}}(-1)=P_G(-1)-1.
\end{equation}
Moreover, \(\alpha(\widehat{G})=\alpha(G)\), because any independent set of
\(\widehat{G}\) containing \(z\) has size \(1\), while every nonempty graph \(G\)
has an independent set of size at least \(1\). Next, we now apply \Cref{cor:PGstar-criterion}.


(a) Let \(G=C_n\), where \(n\ge 3\). Then
 \(\alpha(\widehat{C_n})=\alpha(C_n)=\left\lfloor \frac n2\right\rfloor\). By \Cref{cor:PGstar-criterion}, \(\widehat{C_n}\) is pseudo-Gorenstein\(^{*}\) if and only if
\(P_{\widehat{C_n}}(-1)=(-1)^{\lfloor n/2\rfloor}
\). Using \eqref{eq:cone-eval-minus-one}, this is equivalent to
\[
P_{C_n}(-1)-1=(-1)^{\lfloor n/2\rfloor}.
\]
By \Cref{lem:cycle-minus-one}, we have the following
\[
P_{C_n}(-1)-1=
\begin{cases}
1,& n\equiv 0 \pmod 6,\\[4pt]
0,& n\equiv 1,5 \pmod 6,\\[4pt]
-2,& n\equiv 2,4 \pmod 6,\\[4pt]
-3,& n\equiv 3 \pmod 6.
\end{cases}
\]
Since the right-hand side \( (-1)^{\lfloor n/2\rfloor}\) is always \(\pm 1\), the
only possible case is \(n\equiv 0\pmod 6\), where we must have
\[
1=(-1)^{\lfloor n/2\rfloor}.
\]
If \(n\equiv 0\pmod 6\), then \(\lfloor n/2\rfloor=n/2\), so this occurs exactly
when \(n/2\) is even, i.e. when \(n\equiv 0\pmod{12}\). Hence
\(\widehat{C_n}\) is pseudo-Gorenstein\(^{*}\) if and only if \(n\equiv 0\pmod{12}\).

\smallskip
\noindent
\textbf{(b)} Let \(G=P_n\), where \(n\ge 1\). Then
 \(\alpha(\widehat{P_n})=\alpha(P_n)=\left\lceil \frac n2\right\rceil\).
Again by \Cref{cor:PGstar-criterion}, \(\widehat{P_n}\) is pseudo-Gorenstein\(^{*}\)
if and only if \(P_{\widehat{P_n}}(-1)=(-1)^{\lceil n/2\rceil}\). Using \eqref{eq:cone-eval-minus-one}, this is equivalent to
\[
P_{P_n}(-1)-1=(-1)^{\lceil n/2\rceil}.
\]
By \Cref{lem:path-minus-one}, we have
\[
P_{P_n}(-1)-1=
\begin{cases}
0,& n\equiv 0,5 \pmod 6,\\[4pt]
-1,& n\equiv 1,4 \pmod 6,\\[4pt]
-2,& n\equiv 2,3 \pmod 6.
\end{cases}
\]
Thus equality with \( (-1)^{\lceil n/2\rceil}\in\{\pm1\}\) is possible only when
\(n\equiv 1,4\pmod 6\), in which case we need
\[
-1=(-1)^{\lceil n/2\rceil}.
\]
If \(n\equiv 1\pmod 6\), then this happens exactly for \(n\equiv 1\pmod{12}\).
If \(n\equiv 4\pmod 6\), then this happens exactly for \(n\equiv 10\pmod{12}\).
Therefore \(\widehat{P_n}\) is pseudo-Gorenstein\(^{*}\) if and only if \(n\equiv 1,10 \pmod{12}\).
\end{proof}

\end{document}
